\section{Limitations}
\label{sec:limitations}

This paper is deliberately scoped as an image-level, internal development study on one site-unlabelled public dataset release, and the limitations below follow directly from that scope.

\textit{(1)~No clinician-annotated localization evidence.} All anatomical plausibility judgments in Section~\ref{sec:discussion}.1 are the authors', not an independent radiologist's, and a rating protocol using the existing per-class gallery is planned but awaits clinician participation.

\textit{(2)~Single seed for per-class and XAI-battery results.} Top-line accuracy is three-seed replicated, but the per-class breakdown, consistency results, and rest of the quantitative-XAI battery still come from a single primary-model seed.

\textit{(2b)~The cross-replicate consistency design does not isolate stochastic-seed variance alone.} The three Stage~1 replicate models (Section~\ref{sec:methods}.3) were trained on disjoint subsets of the corrected training set (2,749, 2,089, and 2,661 images respectively), each with a different per-class composition, rather than on identical training data with only the random seed varied. The reported Stage~1 consistency scores $O_r^c$ (Tables~\ref{tbl:5}--\ref{tbl:6}) therefore reflect a combination of re-initialisation, training-volume, and class-composition effects, not isolated stochastic reproducibility. Section~\ref{sec:stage2} reports the cleaner same-data, seed-only version of this experiment for one representative method (Grad-CAM), confirming consistency clears the null-control floor by roughly 2.5$\times$; Section~\ref{sec:stage2-full-battery} then extends this to the full five-method comparison on identical Stage~2 training data, confirming Tables~\ref{tbl:5}--\ref{tbl:6}'s method-by-method ranking holds under the same-data design as well.

\textit{(3)~Bounded perturbation set.} The faithfulness and stability tests use a specific baseline and transform set. Other choices could shift the absolute numbers, though the ranking against the random control is robust.

\textit{(3b)~Uncertainty quantification for XAI-method comparisons.} Tables~\ref{tbl:5}, \ref{tbl:7}, and \ref{tbl:8} now report 2,000-resample image-level bootstrap 95\% confidence intervals alongside every consistency, faithfulness, and stability point estimate, resolving the earlier gap relative to the classification-accuracy results of Section~\ref{sec:results}.1. These intervals confirm what the paper's qualitative caution already suggested: the close faithfulness ranking in Table~\ref{tbl:7} (Grad-CAM 0.480, Grad-CAM++ 0.473, Score-CAM 0.456) has substantially overlapping intervals and should not be over-interpreted, while Grad-CAM++'s and Score-CAM's consistency advantage in Table~\ref{tbl:5} is supported by non-overlapping intervals against the other three methods. Table~\ref{tbl:6}'s per-class breakdown and Table~\ref{tbl:8}'s Pearson/top-10\% IoU columns still report point estimates without per-class or per-metric intervals.

\textit{(3c)~No null control existed for the consistency metric until Stage~2, and it was checked for one method before being extended to all five.} Cosine similarity between non-negative, min-max-normalised heatmaps is structurally biased toward high values, so the original Stage~1 consistency scores (Tables~\ref{tbl:5}--\ref{tbl:6}) had no baseline confirming they exceeded what an untrained model or a synthetic random heatmap pair would already produce. Section~\ref{sec:stage2} closes this gap for Grad-CAM specifically, finding the corrected same-data consistency score clears both an untrained-model floor and a synthetic-noise floor by roughly 2.5$\times$. The null-control comparison itself (untrained-model and synthetic-heatmap floors) has not been recomputed separately for Grad-CAM++, Score-CAM, Saliency, or Eigen-CAM, though Section~\ref{sec:stage2-full-battery}'s full re-run confirms all five methods' Stage~2 consistency scores remain well above where Grad-CAM's floor sits.

\textit{(4)~The margin-based shortcut audit's crop-ablation leg is inconclusive, and a separate batch-correlated shortcut is now confirmed across all nine classes, not merely suspected in one.} The large accuracy drop under margin cropping (0.930~$\rightarrow$~0.712) is confounded with genuine anatomy loss, and that audit tests only for margin-tied shortcuts. Section~\ref{sec:followup-audits} closes this gap in two stages. First, a re-verified, leakage-controlled batch-ID diagnostic still reaches 82.88\% accuracy (8.3$\times$ chance) at distinguishing acquisition batches within one disease class (Carcinoma), and a linear probe on the disease classifier's own \texttt{layer4} features recovers batch identity at 74.75\% (4.5$\times$ chance) for that same class, directly implicating the classifier this paper's XAI battery evaluates. Second, that linear probe is then extended to all nine classes and every eligible batch within each: every class reaches 100.0\% batch-identification accuracy (macro mean 24.0$\times$ chance), confirming the signature is not a Carcinoma-specific artefact. This full-coverage probe uses the same image-level (not cluster-aware) per-batch split as the original single-class probe, so, as discussed where it is reported (Section~\ref{sec:followup-audits}), some of the 100\% ceiling plausibly reflects the same near-duplicate-frame leakage the batch-ID diagnostic's own cluster-aware re-verification found responsible for part of its original inflated figure; a cluster-aware re-split of the nine-class probe has not yet been run and is the natural next step. Neither test rules out shortcuts unrelated to batch identity (scanner-specific artefacts, compression artefacts, or other background cues).

\textit{(4b)~The pHash clustering threshold's sensitivity is now fully characterised, including real accuracy per threshold.} Section~\ref{sec:followup-audits} reports how cluster count, mixed-class-cluster count, and now real Stage~2 test accuracy all move across five threshold values (4 through 12 bits): accuracy rises from 68.44\% at threshold~4 to 97.62\% at threshold~12, a direct demonstration that looser thresholds let leakage back in rather than only a proxy signal (more mixed-class clusters) suggesting it might. This supports the paper's choice of threshold~4 as the more conservative option on every count measured, not only cluster-count and mixed-class proxies. Each threshold's accuracy figure comes from one 10-epoch, single-seed run, lighter than the main Stage~2 retrain's early-stopped, up-to-20-epoch, 4-seed design, so these numbers describe the threshold's effect on accuracy, not a publication-grade model at each threshold.

\textit{(4c)~The filename-prefix-as-acquisition-session inference now has direct, non-circular evidentiary support, though the original dataset providers have not been contacted.} Section~\ref{sec:followup-audits} shows that images sharing a filename prefix cluster tightly in embedded EXIF capture time (median 5.7 minutes), and that no two prefixes occupying the same calendar day ever have overlapping time windows (0 of 3,312 same-day pairs), evidence independent of the filename pattern itself. Thirteen prefixes (5.9\%) span over an hour, three markedly so; these are reported rather than excluded and are not otherwise explained. This resolves the specific circularity objection raised against the filename-only argument, but it does not substitute for confirming the actual export/acquisition convention with Turki et~al., which remains an open item.

\textit{(5)~Site and scanner provenance is unlabelled, not absent.} The source publication states the underlying images were collected across three to four hospitals in Baghdad, Iraq, on four scanner models, but the public release carries no hospital, patient, or scanner identifier. This paper can therefore neither verify the split is balanced across sites and scanners, nor stratify results by site, and no claim of externally validated multicentre or cross-scanner generalizability is made or implied.

\textit{(6)~Noise robustness at the classifier level is now quantified directly, and remains a severe, unresolved weakness; the sweep's own clean-accuracy baseline is itself likely train-set-contaminated.} Section~\ref{sec:followup-audits} reports accuracy under a Gaussian-noise sweep, on the Stage~2 primary checkpoint: accuracy on a stratified subset falls from 90.0\% clean to 20.6\%, close to the 11.1\% chance level, at a noise level ($\sigma=0.12$) that remains visually subtle. That 90.0\% clean-accuracy figure is well above the checkpoint's genuine 70.30\% Stage~2 test accuracy, because the 180-image sample was drawn from the combined Stage~1 training/validation pool rather than the Stage~2 test split, so it likely includes images the model was trained on. The qualitative collapse under noise, applied uniformly across an identical sample at every $\sigma$, is a substantially more consequential finding than any result in the XAI comparison itself, but the 90.0\% starting point specifically should not be quoted as a held-out accuracy figure until the sample is redrawn from the genuine test set. This paper does not attempt to fix the underlying noise sensitivity, only to measure and report it, including this sampling flaw, honestly.

\subsection{Future Research Directions}
\label{sec:future-research}

The remaining steps, in priority order: (1)~Collect clinician gallbladder and pathology region-of-interest annotations and compute IoU, Dice, pointing-game accuracy, and energy-in-ROI for each XAI method. (2)~Run the planned radiologist Likert-rating study and report Cohen's $\kappa$. (3)~Extend the per-class and consistency results to multiple seeds (Limitation~2). (4)~Obtain hospital-, scanner-, and patient-level metadata for the existing release, or prospectively collect site- and patient-labelled data, to enable cross-site and cross-scanner stratification (Limitation~5). (5)~Re-run the nine-class batch-identification linear probe (Section~\ref{sec:followup-audits}) with a cluster-aware, rather than image-level, per-batch split, to establish how much of its 100.0\% per-class ceiling reflects near-duplicate-frame leakage within the probe's own train/validation split versus a genuine, fully leakage-free batch signature (Limitation~4). (6)~Redraw the noise-robustness sweep's 180-image stratified sample from the genuine Stage~2 test split rather than the combined Stage~1 training/validation pool, so its clean-accuracy baseline is not itself likely train-contaminated (Limitation~6). Only once data with verified, site-labelled multicentre provenance become available \cite{ref50}, would a cross-scanner validation study become an appropriate next paper. Throughout, this line of work is intended to follow the FUTURE-AI \cite{ref51}, CLAIM \cite{ref52}, \cite{ref53}, and TRIPOD+AI \cite{ref54} reporting guidelines, and to explicitly audit the generalization pitfalls documented elsewhere in medical-imaging machine learning \cite{ref55}, \cite{ref56}.
